\documentclass[12pt]{article}

% --- Packages ---
\usepackage[utf8]{inputenc}
\usepackage{geometry}
\usepackage{amsmath, amssymb}
\usepackage{hyperref}
\usepackage{graphicx} % Required for \includegraphics
\usepackage{fancyhdr}
\usepackage{enumitem}
\usepackage{autobreak}
\usepackage{amsmath}
\usepackage{breqn}

\usepackage{amsmath,amssymb}
\usepackage{parskip}       % Correct paragraph spacing
\usepackage{xcolor}        % Required for \textcolor command
\usepackage{enumitem}
\usepackage{float}
\usepackage{listings}
\usepackage{xcolor}
\usepackage{tikz}
\usepackage[utf8]{inputenc}
\usepackage[T1]{fontenc}
\usepackage{url}
\usepackage{amsmath,amssymb,amsthm}
\usepackage{booktabs}
\usepackage{pgfplotstable}
\usepackage{siunitx}
\usepackage{geometry}

\setlist{nosep}           % Reduce vertical spacing in lists
\usepackage{amsmath,amssymb}
\usepackage{geometry}
\usepackage{listings}
\usepackage{booktabs}
\usepackage{graphicx}
\usepackage{listings}
\usepackage{pgfplots}
\pgfplotsset{compat=1.18} % sau ultima versiune instalată la tine
\usepackage{filecontents} % dacă 
% --- Page Setup ---
\geometry{a4paper, margin=1in, top=1.2in, headheight=35pt} % Expanded top margin & headheight for the ornament

\lstset{
    language=python,
    basicstyle=\ttfamily\small,
    backgroundcolor=\color{blue!5},
    keywordstyle=\color{purple!80!blue},
    commentstyle=\color{green!50!black},
    stringstyle=\color{red!70!black},
    numberstyle=\tiny\color{gray},
    numbers=left,
    numbersep=15pt,
    frame=shadowbox,
    rulesepcolor=\color{blue!20},
    frameround=tttt,
    breaklines=true,
    captionpos=b,
    showstringspaces=false,
    tabsize=4,
    xleftmargin=20pt,
    xrightmargin=5pt
}

\usepackage[table]{xcolor}
\definecolor{truegreen}{RGB}{144,238,144}
\definecolor{falsered}{RGB}{250,128,114}
\newcommand{\T}{\cellcolor{truegreen}T}
\newcommand{\F}{\cellcolor{falsered}F}

% --- Header & Footer Configuration ---
\pagestyle{fancy}
\fancyhf{} % Clear default headers/footers

% Center the ornament in the header (antet)
\chead{\includegraphics[width=4cm]{ornament_6.pdf}}

% Center the page number in the footer
\cfoot{\thepage}

% Optional: Remove the default rule line below the header
\renewcommand{\headrulewidth}{0pt}

\begin{document}

\begin{titlepage}
    \thispagestyle{fancy} % Force header/footer onto the title page
    \centering
    
    \vspace*{4cm}

    % Main Title
    {\Huge \bfseries Homework 3}

    \vspace{4cm}

    % Name and Date
    {\Large \textbf{Name:} Gatman Vitalie} \\[0.5em]
    {\Large \textbf{Date:} \today} \\[0.5em]
    {\Large \textbf{Group:} FAF-263}

    \vfill
\end{titlepage}

\newpage

% Page 2 content...
\section*{Problem 3.1}
\begin{proof}
\[S(n) \equiv  \sum_{k=1}^{n}{k^2}=\frac{n(n+1)(2n+1)}{6}, \;\;\;\; \forall n \in \mathbb{N}, n \ge 1\]
Let us define $C(n)$ as 
\[C(n)=\left\{n \in \mathbb{N}\;| \;\;\; \sum_{k=1}^{n}{k^2}\not=\frac{n(n+1)(2n+1)}{6}\right\}\]
We asume $C(n)$ is not null so there should exist a least element n.\\
We check $n=1$
\[
\begin{cases}
    {
\sum_{k=1}^{1}{k^2}=1^2=1}\\
    {\frac{1(1+1)(2*1+1)}{6}=\frac{6}{6}=1}
    

\end{cases}
\implies 1 \not\in C(n) \implies n>1
\]
There should exist a least element $n>1$ in $C(n)$ so $n-1 \not\in C(n)$ \\
We prove that for some $n\in N^*$ the $S(n-1)+n^2=S(n)$ does not hold
\[
\begin{cases}
    S(n-1) + n^2 = \sum_{k=1}^{n-1}{k^2}=\frac{(n-1)(n)(2(n-1)+1)}{6} + n^2 = \frac{2n^3-3n^2+n}{6}+n^2 = \frac{2n^3+3n^2+n}{6}\\
    S(n)=\sum_{k=1}^{n}{k^2}=\frac{n(n+1)(2n+1)}{6} = \frac{2n^3+3n^2+n}{6}
\end{cases}
\implies\]\[
\implies S(n-1)+n^2=S(n)
\]
So it holds for all $n\in \mathbb{N}^*$ wich means $C = \varnothing$ and by contradiction $S(n)$ holds for all $n \in \mathbb{N}^*$
\end{proof}


\section*{Problem 3.2}
\begin{enumerate}[label=\alph*)]
  \item\begin{proof}
  Suppose, for contradiction, that the store owner is wrong, so \$12 can be made
  from \$5 and \$8 cards. Define
  \[
    S=\{\, b\in\mathbb{N} \mid \exists\, a\in\mathbb{N} \text{ such that } 5a+8b=12 \,\}.
  \]
  This is the set of possible numbers of \$8 cards used in some way of making \$12.
  By assumption $S$ is nonempty, and $S\subseteq\mathbb{N}$. By the Well Ordering
  Principle, $S$ has a least element; call it $b_0$.

  Becuase $b_0\in S$, there exists $a_0\in\mathbb{N}$ with
  \[
    5a_0+8b_0=12.
  \]
  We know that $5a_0\ge 0$, so $8b_0\le 12$ and $b_0\le 1$. Hence $b_0=0$
  or $b_0=1$.

  \begin{itemize}
    \item If $b_0=0$, then $5a_0=12$, so $a_0=\tfrac{12}{5}\notin\mathbb{N}$,
      a contradiction.
    \item If $b_0=1$, then $5a_0=12-8=4$, so $a_0=\tfrac{4}{5}\notin\mathbb{N}$,
      a contradiction.
  \end{itemize}

  Both cases are impossible, so $S$ must be empty, contradicting our assumption.
  Therefore \$12 cannot be made using only \$5 and \$8 cards, and the store
  owner is correct.
\end{proof}

    \item \begin{proof}
Suppose, for contradiction, that some integer $n \ge 28$ is not formable. Then the set
\[
  C = \{\, n \in \mathbb{N} : n \ge 28 \text{ and } n \text{ is not formable} \,\}
\]
is a nonempty set of natural numbers. By the Well Ordering Principle, $C$ has a least element $m$. Since $m \ge 28$, we have two cases.

\medskip
\noindent\textbf{Case 1: $28 \le m \le 32$.} All elements are formatable:
\begin{align*}
  28 &= 4\cdot 5 + 1\cdot 8, \\
  29 &= 1\cdot 5 + 3\cdot 8, \\
  30 &= 6\cdot 5 + 0\cdot 8, \\
  31 &= 3\cdot 5 + 2\cdot 8, \\
  32 &= 0\cdot 5 + 4\cdot 8.
\end{align*}


\medskip

\noindent\textbf{Case 2: $m \ge 33$.} Then $m - 5 \ge 28$. We can write $m - 5 = 5a + 8b\;\;\;\;\;a, b \ge 0$. Then
\[
  m = 5(a+1) + 8b,
\]
so $m$ is formable, contradicting $m \in C$.

\medskip

Both cases lead to a contradiction, so $C$ is empty. Therefore every amount $n \ge 28$ can be formed using \$5 and \$8 cards.
\end{proof}

\end{enumerate}

\section*{Problem 3.3}


\[
1 + \frac{1}{4} + \frac{1}{9} + \frac{1}{16} + \dots + \frac{1}{n^2} < 2 - \frac{1}{n}, \quad \forall n \in \mathbb{N}, \, n > 1.
\]

\begin{proof}
Lets us define: $\displaystyle f(n)= \sum_{i=1}^{n} \frac{1}{i^2}$.

\paragraph{1. Base case ($n = 2$):}
For $n = 2$, if we evaluate both sides we get:
\[
\text{LHS} = f(2) = 1 + \frac{1}{4} = \frac{5}{4} = 1.25, \quad \text{RHS} = 2 - \frac{1}{2} = \frac{3}{2} = 1.5.
\]
Since $1.25 < 1.5$, the inequality holds true.

\paragraph{2. We prove the induction}
\[
1 + \frac{1}{4} + \frac{1}{9} + \dots + \frac{1}{k^2} < 2 - \frac{1}{k}.
\]
\[
1 + \frac{1}{4} + \dots + \frac{1}{k^2} + \frac{1}{(k+1)^2} < 2 - \frac{1}{k} + \frac{1}{(k+1)^2}.
\]
\[
2 - \frac{1}{k} + \frac{1}{(k+1)^2} \le 2 - \frac{1}{k+1} \iff \frac{1}{(k+1)^2} \le \frac{1}{k} - \frac{1}{k+1}.
\]
\[
\frac{1}{k} - \frac{1}{k+1} = \frac{(k+1) - k}{k(k+1)} = \frac{1}{k(k+1)}.
\]
\[
\frac{1}{(k+1)^2} < \frac{1}{k(k+1)} = \frac{1}{k} - \frac{1}{k+1}.
\]
\[
1 + \frac{1}{4} + \dots + \frac{1}{(k+1)^2} < 2 - \frac{1}{k} + \left(\frac{1}{k} - \frac{1}{k+1}\right) = 2 - \frac{1}{k+1}.
\]
\end{proof}
\section*{Problem 3.4}
\[
F_1 + F_2 + F_3 + \dots + F_n = F_{n+2} - 1, \quad \forall n \in \mathbb{N}, \, n \ge 1.
\]

\begin{proof}
From the definition of the Fibonacci sequence: $F_1 = 1$, $F_2 = 1$, and $F_k = F_{k-1} + F_{k-2}$ for all integers $k \ge 3$. 

We define be the statement: $\displaystyle f(n) = \sum_{i=1}^{n} F_i = F_{n+2} - 1$.

\paragraph{1. Base Case ($n = 1$):}
For $n = 1$, if we evaluate both sides:
\[
\text{LHS} = f(1) = F_1 = 1, \quad \text{RHS} = F_{1+2} - 1 = F_3 - 1 = 2 - 1 = 1.
\]
Since $\text{LHS} = \text{RHS}$, the inequiality holds true.

\paragraph{2. We prove the induction}
\[
F_1 + F_2 + F_3 + \dots + F_k = F_{k+2} - 1.
\]
\[
F_1 + F_2 + \dots + F_k + F_{k+1} = F_{(k+1)+2} - 1 = F_{k+3} - 1.
\]
\begin{align*}
\sum_{i=1}^{k+1} F_i &= \left( \sum_{i=1}^{k} F_i \right) + F_{k+1} \\
&= (F_{k+2} - 1) + F_{k+1}  \\
&= (F_{k+2} + F_{k+1}) - 1.
\end{align*}
We can rewrite $F_{k+3}$ as, $F_{k+3} = F_{k+2} + F_{k+1}$. Substituting this in:
\[
\sum_{i=1}^{k+1} F_i = F_{k+3} - 1.
\]
Thus the equality holds true.
\end{proof}
\newpage
\section*{Problem 3.5}
\begin{proof}
\[4^n+7\le 5^n\]
\paragraph{1. Base Case ($n = 2$):}
For $n = 2$, we evaluate both sides:
\[
\text{LHS} = 4^2 + 7 = 16 + 7 = 23, \quad \text{RHS} = 5^2 = 25.
\]
Since $23 \le 25$, the equality holds true.

\paragraph{2. We prove the induction}
\[
4^k + 7 \le 5^k.
\]
\begin{align*}
4^{k+1} + 7 &= 4 \cdot 4^k + 7 \\
&< 5 \cdot 4^k + 35 \quad \text{(since } 4 \cdot 4^k < 5 \cdot 4^k \text{ and } 7 < 35\text{)} \\
&= 5(4^k + 7).
\end{align*}
\[
5(4^k + 7) \le 5 \cdot (5^k) = 5^{k+1}.
\]
\[
4^{k+1} + 7 \le 5^{k+1}.
\]
Thus, the inequality holds.
\end{proof}
\section*{Problem 3.6}

A positive integer $n$ is called a \textbf{perfect number} if it is equal to the sum of its proper positive divisors.
Prove by induction that if $2^k - 1$ is a prime number (a \textbf{Mersenne prime}), then $n = 2^{k-1}(2^k - 1)$ is an even perfect number.

\begin{proof}
Let $p = 2^k - 1$ be a Mersenne prime. Since $p$ is prime and $p > 1$, we must have $k \ge 2$, so $n = 2^{k-1}p$ is divisible by $2$ and is therefore even.

The proper positive divisors of $n = 2^{k-1}p$ consist of:
\begin{enumerate}
    \item Divisors that are powers of 2: $1, 2, 2^2, \dots, 2^{k-1}$.
    \item Divisors multiplied by $p$ (excluding $n$ itself): $p, 2p, 2^2p, \dots, 2^{k-2}p$.
\end{enumerate}

Thus, the sum of all proper positive divisors of $n$ is:
\[
S = \sum_{i=0}^{k-1} 2^i + p \sum_{i=0}^{k-2} 2^i.
\]

We first prove by induction that for any integer $m \ge 0$:
\[
g(m) = \sum_{i=0}^{m} 2^i = 2^{m+1} - 1.
\]

\paragraph{1. Base Case ($m = 0$):}
For $m = 0$:
\[
\text{LHS} = g(0) = 2^0 = 1, \quad \text{RHS} = 2^{0+1} - 1 = 2 - 1 = 1.
\]

\paragraph{2. We prove the induction}
\[
\sum_{i=0}^{j+1} 2^i = \left(\sum_{i=0}^{j} 2^i\right) + 2^{j+1} = (2^{j+1} - 1) + 2^{j+1} = 2 \cdot 2^{j+1} - 1 = 2^{j+2} - 1.
\]
Thus, the sum formula holds for all $m \ge 0$.

Using $g(k-1)$ and $g(k-2)$:
\[
\sum_{i=0}^{k-1} 2^i = 2^k - 1 = p, \quad \sum_{i=0}^{k-2} 2^i = 2^{k-1} - 1.
\]
\[
S = p + p(2^{k-1} - 1) = p\left(1 + 2^{k-1} - 1\right) = p \cdot 2^{k-1} = n.
\]
Since the sum of the proper positive divisors of $n$ equals $n$, $n$ is a perfect number. Since $k \ge 2$, $n$ is an even perfect number.
\end{proof}


\section*{Problem 3.7}
\begin{proof}
    \[
F_n < 2^n.
\]
\paragraph{1. Base Cases ($n = 1$ and $n = 2$):}
For $n = 1$:
\[
\text{LHS} = F_1 = 1, \quad \text{RHS} = 2^1 = 2.
\]
Since $1 < 2$, the inequality holds true.

\vspace{0.5em}
\noindent For $n = 2$:
\[
\text{LHS} = F_2 = 1, \quad \text{RHS} = 2^2 = 4.
\]
Since $1 < 4$, the inequality holds true.

\paragraph{2. We prove the induction}
Assume by strong induction that $F_i < 2^i$ holds for all $1 \le i \le k$, where $k \ge 2$. We show that $F_{k+1} < 2^{k+1}$.

By definition of the Fibonacci sequence:
\[
F_{k+1} = F_k + F_{k-1}.
\]
Applying the strong inductive hypothesis to $F_k$ and $F_{k-1}$ (since $1 \le k-1 < k$):
\[
F_k < 2^k \quad \text{and} \quad F_{k-1} < 2^{k-1}.
\]
Substituting these into the expression for $F_{k+1}$:
\[
F_{k+1} = F_k + F_{k-1} < 2^k + 2^{k-1}.
\]
Since $2^{k-1} < 2^k$ for all $k \ge 2$:
\[
2^k + 2^{k-1} < 2^k + 2^k = 2 \cdot 2^k = 2^{k+1}.
\]
Therefore:
\[
F_{k+1} < 2^{k+1}.
\]
Thus, the inequality holds.
\end{proof}


\section*{Problem 3.8}
\begin{proof}
Let $P(n)$ be the statement that an amount of $\$$ can be formed using natural number quantities of $5\$$ and $8\$$ cards: $n = 5a + 8b$ for some $a, b \in \mathbb{N}$.

\paragraph{1. Base Cases ($n = 28, \,29, \,30,\, 31,\, 32$):}
We evaluate $P(n)$ for five consecutive base values:
\begin{itemize}
    \item For $n = 28$: $28 = 5(4) + 8(1)$.
    \item For $n = 29$: $29 = 5(1) + 8(3)$.
    \item For $n = 30$: $30 = 5(6) + 8(0)$.
    \item For $n = 31$: $31 = 5(3) + 8(2)$.
    \item For $n = 32$: $32 = 5(0) + 8(4)$.
\end{itemize}
Because all are valide the base case is true.

\paragraph{2. We prove the induction}
We assume by strong induction that $P(i)$ is true for all $28 \le i \le k$, where $k \ge 32$. We prove that $P(k+1)$ is true.

\[
(k+1) - 5 = k - 4 \ge 28.
\]
Since $28 \le k - 4 \le k$, by the strong inductive hypothesis $P(k-4)$ holds true, so there exist $a', b' \in \mathbb{N}$ such that:
\[
k - 4 = 5a' + 8b'.
\]
Adding $\$5$ to both sides yields:
\[
k + 1 = (k - 4) + 5 = (5a' + 8b') + 5 = 5(a' + 1) + 8b'.
\]
Since $a' \ge 0 \implies a' + 1 \in \mathbb{N}$ and $b' \in \mathbb{N}$, this provides a valid construction for $\$(k+1)$ using $\$5$ and $\$8$ cards.

Thus, $P(k+1)$ his true.\end{proof}


\newpage
\section*{Problem 3.9}

\textbf{Flaw in the Proof:}

The flaw in the bogus proof occurs in the \textbf{inductive step} when \(n = 0\).

Specifically, the proof relies on the algebraic identity:
\[
a^{n+1} = \frac{a^n \cdot a^n}{a^{n-1}}
\]
and substitutes \(a^n = 1\) and \(a^{n-1} = 1\) using the induction hypothesis\\ \(P(n) : \forall k \le n, \, a^k = 1\).

\paragraph{1. Invalid Exponent for \(n = 0\):}
When \(n = 0\), the expression requires \(a^{n-1} = a^{-1}\):
\[
a^{0+1} = a^1 = \frac{a^0 \cdot a^0}{a^{-1}}
\]
The induction hypothesis \(P(0)\) only asserts that \(a^k = 1\) for non-negative integers \(k \le 0\) (i.e., \(k = 0\), giving \(a^0 = 1\)). It provides no information about \(k = -1\), because \(-1\) is not a non-negative integer (\(k \notin \mathbb{N}\)). Thus, substituting \(a^{-1} = 1\) is invalid.

\paragraph{2. Breakdown of the Inductive Chain:}
For \(a^{n-1} = 1\) to follow legitimately from the induction hypothesis \(P(n)\), \(n - 1\) must be a non-negative integer, which requires \(n \ge 1\). Consequently, \(P(0)\) cannot imply \(P(1)\). Since the transition from \(n = 0\) to \(n = 1\) fails, the entire inductive argument collapses.


\section*{Problem 3.10}

\begin{proof}
Let \((x, y)\) denote Romeo's coordinates on the grid. We define the invariant quantity:
\[
I(x, y) = x + 2y \pmod 7.
\]

\paragraph{1. Initial State:}
At Romeo's starting position \((0, 0)\), we evaluate the invariant:
\[
I(0, 0) = 0 + 2(0) = 0 \equiv 0 \pmod 7.
\]

\paragraph{2. Invariance under moves:}
Suppose Romeo is at position \((x, y)\) such that \(x + 2y \equiv 0 \pmod 7\). Let \((x', y')\) be his new position after applying any of the four allowed moves:
\begin{enumerate}
    \item Move 1 \((+2, -1)\): \(x' + 2y' = (x + 2) + 2(y - 1) = x + 2y + 0 \equiv 0 \pmod 7\).
    \item Move 2 \((-2, +1)\): \(x' + 2y' = (x - 2) + 2(y + 1) = x + 2y + 0 \equiv 0 \pmod 7\).
    \item Move 3 \((+1, +3)\): \(x' + 2y' = (x + 1) + 2(y + 3) = x + 2y + 7 \equiv 0 \pmod 7\).
    \item Move 4 \((-1, -3)\): \(x' + 2y' = (x - 1) + 2(y - 3) = x + 2y - 7 \equiv 0 \pmod 7\).
\end{enumerate}
In every case, \(x' + 2y' \equiv 0 \pmod 7\). By induction on the number of moves, any position reachable by Romeo must satisfy \(x + 2y \equiv 0 \pmod 7\).

\paragraph{3. Target Evaluation:}
For Julietta's position at \((1, 1)\), we have:
\[
1 + 2(1) = 3 \not\equiv 0 \pmod 7.
\]
Since \(3 \not\equiv 0 \pmod 7\), the point \((1, 1)\) does not satisfy the invariant. Therefore, Romeo will never be able to reach Julietta.
\end{proof}



\section*{Problem 3.11}
\begin{proof}
We write $(B, G, K)$ for the numbers of blue, green, and black chameleons,
respectively, so the initial state is $(11, 13, 17)$.
 
\paragraph{1. Invariant.}
In any meeting of two chameleons of different colors, the two colors that
meet each decrease by $1$ and the third color increases by $2$. For example,
\[
  \text{Blue \& Green meet:} \quad (B, G, K) \longmapsto (B-1,\; G-1,\; K+2).
\]
Thus each of $B, G, K$ changes by either $-1$ or $+2$ in every meeting. Since
$-1 \equiv 2 \pmod 3$, all three counts change by the same residue modulo $3$.
Therefore each of the differences
\[
  B - G, \qquad G - K, \qquad K - B
\]
is unchanged modulo $3$ by every meeting. Explicitly, in the example above,
\begin{align*}
  (B-1) - (G-1) &= B - G, \\
  (G-1) - (K+2) &= (G - K) - 3 \equiv G - K \pmod 3, \\
  (K+2) - (B-1) &= (K - B) + 3 \equiv K - B \pmod 3,
\end{align*}
and the other two types of meeting are identical up to relabeling colors.
 
\paragraph{2. All blue and all black are impossible.}
Initially,
\[
  B - G = 11 - 13 = -2 \equiv 1, \qquad
  G - K = 13 - 17 = -4 \equiv 2, \qquad
  K - B = 17 - 11 = 6 \equiv 0 \pmod 3.
\]
\begin{itemize}
  \item \textbf{All blue}, $(41, 0, 0)$: here $G - K = 0$. By the invariant we
  would need $G - K \equiv 0 \pmod 3$ initially, but $G - K \equiv 2 \pmod 3$.
  Impossible.
  \item \textbf{All black}, $(0, 0, 41)$: here $B - G = 0$. By the invariant we
  would need $B - G \equiv 0 \pmod 3$ initially, but $B - G \equiv 1 \pmod 3$.
  Impossible.
\end{itemize}
 
\paragraph{3. All green is achievable.}
The invariant does not rule out $(0, 41, 0)$, since $K - B \equiv 0 \pmod 3$
initially. In fact it can be reached:
\begin{enumerate}
  \item Let a black and a green chameleon meet, twice in succession:
  \[
    (11, 13, 17) \longrightarrow (13, 12, 16) \longrightarrow (15, 11, 15).
  \]
  \item Now $B = K = 15$. Let a blue and a black chameleon meet $15$ times.
  Each meeting changes $(B, G, K)$ by $(-1, +2, -1)$, so the result is
  \[
    (15 - 15,\; 11 + 2 \cdot 15,\; 15 - 15) = (0, 41, 0).
  \]
\end{enumerate}
 
Hence all blue and all black are impossible, while all green is achievable;
in particular, the statement that no monochromatic state can be reached is
false for the initial state $(11, 13, 17)$.
\end{proof}
 

\section*{Problem 3.12}

\begin{proof}

We analyze the number of black and white squares on the chessboard.

\paragraph{1. Color Invariant:}

Color the \(8\times8\) chessboard in the usual alternating pattern of
black and white squares. Since the chessboard has an even number of
squares in each row and column, opposite corners have the same color.

Initially, the chessboard contains
\[
32 \text{ black squares} \qquad\text{and}\qquad 32 \text{ white squares}.
\]

Removing two opposite corners therefore removes two squares of the same
color. Without loss of generality, suppose both removed squares are
black. The remaining board then contains
\[
30 \text{ black squares} \qquad\text{and}\qquad 32 \text{ white squares}.
\]

\paragraph{2. Effect of a Domino:}

A domino always covers two adjacent squares. Since adjacent squares on a
chessboard have opposite colors, every domino covers exactly one black
square and one white square.

Therefore, if \(31\) dominoes were used, they would cover exactly
\[
31 \text{ black squares} \qquad\text{and}\qquad
31 \text{ white squares}.
\]

\paragraph{3. Contradiction:}

However, after removing the two opposite corners, the chessboard
contains \(30\) black squares and \(32\) white squares. Thus, the
numbers of black and white squares are not equal.

Since every domino covers one square of each color, it is impossible for
\(31\) dominoes to cover the remaining board.

Therefore the chessboard cannot be tiled with 31 dominoes.

\end{proof}


\section*{Problem 3.13}
\begin{proof}

We define by \(n_L(w)\) the number of left brackets and by \(n_R(w)\) the number of right brackets in \(w\). We wish to prove that for all \(w \in BM\), \(n_L(w) = n_R(w)\).

\paragraph{Base Case}
For \(w = \lambda\), \(n_L(\lambda) = 0\) and \(n_R(\lambda) = 0\).
Since \(0 = 0\), \(n_L(\lambda) = n_R(\lambda)\) holds.

\paragraph{We prove the induction}
Let \(w \in BM\) be generated by the constructor rule, so \(w = [s]t\) for some \(s, t \in BM\).\\
We asume \(n_L(s) = n_R(s)\) and \(n_L(t) = n_R(t)\).\\
By counting brackets in \(w = [s]t\):\\
\begin{align*}
n_L(w) &= 1 + n_L(s) + n_L(t) \\
n_R(w) &= 1 + n_R(s) + n_R(t)
\end{align*}
Applying the induction hypothesis \(n_L(s) = n_R(s)\) and \(n_L(t) = n_R(t)\):
\[
n_L(w) = 1 + n_L(s) + n_L(t) = 1 + n_R(s) + n_R(t) = n_R(w).
\]
Thus, \(n_L(w) = n_R(w)\) holds for \(w = [s]t\).
\end{proof}

\section*{Problem 3.14}
\begin{enumerate}[label=\alph*)]
    \item $1/x$
    \begin{enumerate}[label=\arabic*.]
        \item $A=x$
        \item $B=2^A=2^x$
        \item $C=B^{-1}=\log_2{x}$
        \item $D=2^{-1\times C} = 2^{-\log_{2}{x}}=\left(2^{\log_2{x}}\right)^{-1}=x^{-1}=\frac{1}{x}$
    \end{enumerate}
    \item We prove that $\tan(x)$ is in EF, them that $\frac{x-5}{x+5}$ is in EF 
    \begin{enumerate}[label=\arabic*.]
        \item $A=x$
        \item $B=\frac{1}{x}$
        \item $C=\sin(x)$
        \item $D=AA=xx=x^2$
        \item $E=C\circ\left(x+\frac{\pi}{2}\right)=\cos{x}$
        \item $F=C \times B \circ E = \sin{x}\times\frac{1}{\cos{x}}=\tan{x}$
        \item $G=(A-5) \times B \circ (A+5)=\frac{x-5}{x+5}$
        \item $H=F \circ G = \tan {\left(\frac{x-5}{x+5}\right)}$
    \end{enumerate}
    \item $\ln{\left(x^2+1\right)}$
    \begin{enumerate}[label=\arabic*.]
        \item $A=x$
        \item $B=\frac{1}{x}$
        \item $C=\log_2{x}$
        \item $D=AA=xx=x^2$ where $e=2.71\cdots$ and is the euler number
        \item $E=C\circ A \times B\circ (C\circ e)=\log_2{x}\times \frac{1}{\log_2{e}}=\frac{\log_2{x}}{\log_2{e}}=\ln{x}$
        \item $F=A\times A + 1 = x^2 + 1$
        \item $H=E \circ F = \ln{\left(x^2+1\right)}$
    \end{enumerate}
    \item Proof by structural induction that if \(f \in EF\), then \(f' \in EF\).
    \begin{enumerate}[label=\arabic*.]
        \item \textbf{Base Cases:}
        \begin{enumerate}[label=\roman*.]
            \item \(f(x) = x \implies f'(x) = 1 \in EF\) (constant)
            \item \(f(x) = c \implies f'(x) = 0 \in EF\) (constant)
            \item \(f(x) = \sin(x) \implies f'(x) = \cos(x) \in EF\) (proven in part b)
        \end{enumerate}
        \item \textbf{Inductive Hypothesis:} WE assume \(f, g \in EF\) with \(f', g' \in EF\).
        \item \textbf{Inductive Step:}
        \begin{enumerate}[label=\roman*.]
            \item \((f + g)' = f' + g' \in EF\) (sum rule)
            \item \((fg)' = f' \times g + f \times g' \in EF\) (product rule)
            \item \((2^g)' = 2^g \times g' \times \ln(2) \in EF\) (\(\ln(2) = \frac{1}{\log_2{e}} \in EF\) because is constant)
            \item \((f^{-1})' = B \circ f' \circ f^{-1} \in EF\)  \(B(x) = \frac{1}{x} \in EF\) (provem in part a)
            \item \((f \circ g)' = (f' \circ g) \times g' \in EF\) (chain rule)
        \end{enumerate}
    \end{enumerate}
\end{enumerate}
\begin{center}
\fontsize{11}{9}\selectfont
Vitalie is NS (natural stupidity) and can make mistakes. Please perform two factor verification on the Domashneye Rabota
\end{center}
\end{document}